\documentclass[12pt]{article}
\usepackage{fullpage}
\usepackage{enumitem}
\usepackage{listings,amssymb}
\usepackage{listings}
\usepackage{hyperref}
\usepackage{graphicx,xcolor}
\usepackage[T1]{fontenc}
\usepackage{newtxtext,newtxmath}

\lstset{numbers=left, numberstyle=\tiny}
% TODO: Reconfirm HW5 assignment point total before publication.
% Current split assumes the present 70-point Gradescope placeholder (35/35).

\title{CS 147: Lab 05 \\ Due Date: November 12, 2026}
\author{}
\date{\today}

\begin{document}
\frenchspacing
\maketitle

\section*{Homework: Cache Controller Design}

\subsection*{Assignment Overview}

Read Section 5.9 of the H\&P textbook before starting.  You may also
want to read Appendix B-10 of the textbook for a detailed description
of Moore/Mealy state machines.

The goal of this assignment is to develop:

\begin{itemize}
\item A schematic of the memory system
\item A state machine diagram for the cache controller
\end{itemize}

These designs correspond to the cache system that will be implemented later in the project.
All of the cache details are listed in Project Phase 2 writeup (see Section for Phase \# 2.3).
\textbf{Read it carefully before starting!}.

As explained there, you will need to determine how your cache is arranged and functions before starting
implementation. 
Before beginning implementation, you should determine how your cache is arranged and how it functions.
Drawing the state machine for your cache controller will help ensure that the controller logic behaves correctly.

You may implement either a \textbf{Mealy} or \textbf{Moore} state machine, although a Moore machine is recommended since it is generally easier to reason about for this type of controller.

Be aware that the resulting state machine may be fairly large, so starting early is recommended.

\subsection*{Problem 1: Direct-Mapped Cache (35 points)}

Create the schematic and state machine diagram for the \textbf{direct-mapped cache} that will be used in the next phase of the project.

\paragraph{Memory System Schematic}

Your schematic should show how the following components connect together:

\begin{itemize}
\item Direct-mapped cache
\item Cache controller
\item Four-bank memory module
\end{itemize}

These are the same modules that will later be connected together inside the memory system module used in the project.
(\textsl{NOTE: these are the modules you will
eventually use inside \texttt{mem\_system.v} in Phase 2.3}). And the state machine diagram should show the various
states necessary to correctly operate your cache and the transitions between those states.

\paragraph{Cache Controller State Machine}

Create a state machine diagram for the cache controller.
The diagram should show:

\begin{itemize}
\item The states required to correctly operate the cache
\item The transitions between those states
\item The conditions that trigger each transition
\end{itemize}

\subsection*{Problem 2: Two-Way Set Associative Cache (35 points)}

Next, modify your design to support a \textbf{two-way set associative cache}.

Create the following:

\begin{itemize}
\item A schematic of the memory system showing the two-way set associative cache, cache controller, and four-bank memory module
\item A state machine diagram for the cache controller
\end{itemize}

If your state machine from Problem 1 does not require changes, you may submit the same state machine diagram for both problems.

\section*{Files to Submit}

Place the following files in the assignment directories:

\begin{verbatim}
assignments/hw05/
  hw5_1/
    schematic.pdf
    cacheFSM.pdf

  hw5_2/
    schematic.pdf
    cacheFSM.pdf
\end{verbatim}

\begin{itemize}
\item \texttt{hw5\_1/} should contain the files for Problem 1
\item \texttt{hw5\_2/} should contain the files for Problem 2
\item If the same file is used for both problems, place a copy in each directory
\end{itemize}

Digitally created diagrams are strongly preferred.  An easy to use tool for drawing state machines
is \url{https://madebyevan.com/fsm/}.
Hand-drawn diagrams will be accepted only if they are very neat and clearly legible.
If the schematic or the FSM diagram is missing, that component of the problem will receive zero points.

\section*{Submission Instructions}

To build a submission for grading, run:

\begin{verbatim}
./run build_submission hw05
\end{verbatim}

This will package the expected files into a submission zip and place it in:

\begin{verbatim}
generated_turnins/hw05/
\end{verbatim}

Upload the generated zip file to the corresponding Gradescope assignment.

\end{document}
